\section*{Chapitre \ref{ch:m3:power-markets-design} --- Marchés de l'électricité~I~: conception}

\begin{solution}{exo:m3:power-markets-design:1}
\qty{40}{EUR/MWh}~: les 300~MWh à zéro et 150 des 200~MWh à 40 sont acceptés, l'offre
à 40 étant marginale~; les offres à 90 et 150 sont rejetées.
\end{solution}

\begin{solution}{exo:m3:power-markets-design:2}
$(30 + 0.202 \times 70)/0.5 = \qty{88.28}{EUR/MWh}$.
\end{solution}

\begin{solution}{exo:m3:power-markets-design:3}
Dans une adjudication à prix uniforme, toute l'énergie acceptée vaut le même prix à la marge, et
les vendeurs peuvent offrir leur coût marginal~: toute offre acceptée rapporte le prix de compensation.
Avec le paiement au prix offert, chaque vendeur offrirait ce qu'il anticipe comme prix marginal, et non son
coût, de sorte que les offres ne révéleraient plus l'\omterm{def:m3:power-markets-design:merit}{ordre de préséance économique} (merit order)~; le format d'enchère lui-même est
étudié dans le livre~4, chapitre~29.
\end{solution}

\begin{solution}{exo:m3:power-markets-design:4}
Flux de 3~GW de A vers B, prix \qty{82.38}{} et \qty{89.68}{EUR/MWh}, \omterm{def:m3:power-markets-design:lmp}{rente de congestion} $7.30 \times
3\,000 \approx$ 21\,900~EUR pour l'heure (21\,890~EUR avec l'écart de prix non arrondi,
\qty{7.297}{EUR/MWh}). La zone A dispose de 63~GW de capacité moins chère que
\qty{89.35}{} face à 50~GW de demande~: avec plus de 13~GW de ligne, ses centrales au gaz peuvent
exporter également, et les prix convergent à \qty{89.35}{EUR/MWh}.
\end{solution}

\begin{solution}{exo:m3:power-markets-design:5}
Elle produit tant que le prix augmenté de la prime est positif~: elle s'arrête en dessous de
$-\qty{60}{EUR/MWh}$, soit pendant les quatre heures de 12:00 à 16:00 le 11 mai 2025
(prix allant de $-212.82$ à $-110.06$).
\end{solution}

\begin{solution}{exo:m3:power-markets-design:6}
À l'intérieur d'une zone, le réseau présente encore des limites~; le prix nodal, qui intègre la congestion,
varie d'un nœud à l'autre, mais le marché zonal ignore cela et s'équilibre à un seul prix. Le
gestionnaire du réseau procède alors à un redispatching, en payant les centrales situées d'un côté du goulet pour qu'elles baissent leur production et
celles de l'autre côté pour qu'elles l'augmentent, et le coût est répercuté sur
l'ensemble des consommateurs par les tarifs d'acheminement.
\end{solution}

\begin{solution}{exo:m3:power-markets-design:7}
457 en 2024 et 573 en 2025~; le maximum à 13:00 heure locale (171 sur les deux années)~; le
prix le plus bas, $-\qty{250.32}{EUR/MWh}$, à 13:00 le 11 mai 2025.
\end{solution}

\begin{solution}{exo:m3:power-markets-design:8}
Un prix négatif signifie qu'à cette heure un MWh de plus a une valeur négative~: l'offre qui ne peut
ou ne veut pas s'arrêter excède la demande. Il rémunère les consommateurs flexibles pour qu'ils consomment et les producteurs flexibles
pour qu'ils s'arrêtent, ce qui est le marché qui fonctionne~; les producteurs qui continuent de produire le font parce que s'arrêter
coûte plus cher (redémarrages) ou parce qu'ils sont rémunérés au MWh produit (subventions).
\end{solution}

\begin{solution}{pb:m3:power-markets-design:1}
\textbf{1.} Neuf heures, moyenne $-\qty{98.54}{EUR/MWh}$. \textbf{2.} $-\qty{250.32}{}$, 13:00
à 14:00. \textbf{3.} De 11:00 à 16:00 (cinq heures). \textbf{4.} Exportations vers les zones voisines
jusqu'à la capacité des lignes, le stockage (stations de transfert d'énergie par pompage, batteries) et l'écrêtement de
la production des producteurs qui s'arrêtent à prix négatifs. \textbf{5.} L'offre proposée sous zéro, venant de producteurs qui
perdent davantage en s'arrêtant ou sont rémunérés au MWh produit, a excédé la demande à zéro.
\textbf{6.} $400 \times (82.38 + 250.32) = $ 133\,080~EUR. \textbf{7.} $400 \times \sum_h
(82.38 - p_h) = $ 651\,308~EUR sur les neuf heures. \textbf{8.} L'arrêt évite cette
perte mais coûte un redémarrage (combustible pour la remise en chauffe, usure et risque de ne pas être prêt pour la
pointe du soir). \textbf{9.} Rester en service si le coût du redémarrage dépasse la perte subie en fonctionnant au
minimum technique pendant les heures négatives. \textbf{10.} Une offre bloc sur la journée à charge minimale, avec
une limite qui intègre le coût de redémarrage, ou des offres à prix négatifs pour la charge minimale.
\textbf{11.} Chaque heure négative dont le prix est supérieur à $-60$~: 9:00, 10:00, 11:00, 16:00 et
17:00. \textbf{12.} Environ $-\qty{60}{EUR/MWh}$~: leurs offres se situent à moins leur prime.
\textbf{13.} En ne versant aucune prime lorsque les prix sont négatifs~: en Allemagne, pour les nouvelles installations solaires
depuis février 2025, à chaque quart d'heure négatif, de sorte que ces producteurs offrent à zéro ou
s'arrêtent. \textbf{14.} Il achète (il charge) pendant les heures les moins chères, ajoutant ainsi de la demande, et vend
le soir. \textbf{15.} Les consommateurs flexibles, le stockage et les producteurs capables de s'arrêter et de racheter
à bas coût leurs ventes à terme. \textbf{16.} Davantage d'exportations aurait fait monter le prix allemand
et fait baisser celui des voisins, jusqu'à ce que la ligne soit saturée ou que les prix se rejoignent. \textbf{17.} Un fort ensoleillement,
des températures douces et des jours fériés donnent une forte production solaire face à une faible demande. \textbf{18.} Elle est
rémunérée pour consommer~: sa valeur est le prix négatif multiplié par la consommation qu'elle peut déplacer.
\textbf{19.} \emph{Résultat nommé~:} l'unité devrait s'arrêter si le redémarrage coûte moins que
651\,308~EUR. \textbf{20.} Le système dispose d'une offre inflexible supérieure à la demande à cette
heure, et rémunère quiconque peut l'absorber ou s'arrêter.
\end{solution}

\begin{solution}{iq:m3:power-markets-design:1}
Par une enchère quotidienne à offres scellées, à midi, pour chaque \omterm{def:m3:power-markets-design:mtu}{unité de temps de marché} (MTU) du lendemain~: les offres d'achat et de vente forment
des courbes en escalier~; un seul algorithme compense ensemble toutes les zones couplées, en maximisant le bien-être sous
réserve des capacités transfrontalières, et chaque zone obtient un prix uniforme au point où ses courbes se croisent.
\iqlookfor{prix uniforme, maximisation du bien-être, \omterm{def:m3:power-markets-design:coupling}{couplage des marchés}.}
\end{solution}

\begin{solution}{iq:m3:power-markets-design:2}
La pile de l'empilement de l'offre classée par coût marginal. L'éolien a un coût marginal quasi nul, ce qui décale la
courbe vers la droite et la demande la rencontre plus bas~: les prix baissent aux heures venteuses, d'autant plus que
l'empilement est pentu.
\iqlookfor{décalage de l'empilement, cannibalisation des heures venteuses.}
\end{solution}

\begin{solution}{iq:m3:power-markets-design:3}
Prévisions de charge, de vent et de solaire (charge résiduelle), prix du combustible et du carbone, disponibilité des centrales,
capacités des interconnexions et charge résiduelle des voisins, effets de calendrier. Au printemps, les prévisions solaires
et éoliennes dominent les prix de midi.
\iqlookfor{charge résiduelle comme variable clé~; météo.}
\end{solution}

\begin{solution}{iq:m3:power-markets-design:4}
Lorsque l'offre inflexible ou subventionnée dépasse la demande augmentée des exportations~: les producteurs paient pour éviter
les coûts de redémarrage ou pour conserver les subventions. Le minimum harmonisé de l'enchère le borne (en Europe
$-\qty{500}{}$, puis $-\qty{600}{EUR/MWh}$ à partir de mai 2026).
\iqlookfor{mécanismes, et connaissance du plancher.}
\end{solution}

\begin{solution}{iq:m3:power-markets-design:5}
Avec une tarification zonale, il reçoit le prix de la zone quelle que soit la congestion locale, mais peut être
écrêté par le redispatching~; avec une tarification nodale, il reçoit le prix de son nœud, qui baisse lorsque
les lignes locales sont congestionnées, il supporte donc le risque de congestion et a besoin de droits de transport pour
se couvrir.
\iqlookfor{qui supporte le risque de congestion, et la couverture.}
\end{solution}

\begin{solution}{iq:m3:power-markets-design:6}
Un problème de maximisation du bien-être sous contraintes de réseau (un programme linéaire ou quadratique
par scénario) auquel s'ajoute une couche combinatoire pour les blocs (séparation et évaluation, ou une heuristique avec
contrôles de paradoxes). Vitesse~: démarrage à chaud à partir du scénario précédent, décomposition par jour, vectorisation
des courbes par MTU, et mise en cache des décisions de blocs qui ne changent pas.
\iqlookfor{programmation linéaire et combinatoire, et accélération pratique.}
\end{solution}
